%% kapa: exam=1 round=35 year=2024 period=1 form=A subject=economics subject_ko=경제학원론 seq=35 number=75
%% source: Q-Net 감정평가사 기출문제 (cst00309) artlSeq=5213229 file=raw/1008/5213229/1교시 문제지 원본.pdf
%% conversion: kordoc 4.24.51 -> tools/gen1.py (auto); review=proofread-sonnet; flags=none
\begin{question}{75}
폐쇄경제 IS-LM모형에 관한 설명으로 옳은 것은?
\begin{choices}
  \item 화폐수요의 이자율 탄력성이 0이면 경제는 유동성함정(liquidity trap) 상태에 직면한다.
  \item LM곡선이 수직선이고 IS곡선이 우하향하면, 완전한 구축효과(crowding-out effect)가 나타난다.
  \item IS곡선이 수평선이고 LM곡선이 우상향하면, 통화정책은 국민소득에 영향을 미치지 않는다.
  \item 소비가 이자율에 영향을 받을 때, 피구효과(Pigou effect)가 발생한다.
  \item IS곡선이 우하향할 때, IS곡선의 위쪽에 있는 점은 생산물시장이 초과수요 상태이다.
\end{choices}
\end{question}
\answer{2}
